^^A Measurements behind the decisions on page breaks.  \input by appendix.tex, so
^^A no % comments, only ^^A, and no literal vertical bar, as in that file.
^^A The numbers come from measure-real/ and measure-flow/; see measure-real/README.md.

\section{Measurements on real documents}\label{app:pb-measure}

This section reports what the page-break defaults do to real documents.
The numbers are from the package at commit \texttt{de43f08}.
At \texttt{3c059c3}, which adds one \texttt{.aux} write and a whatsit after a lead-in, every page of
all eight documents has the same words, with and without the probe, so the tables hold there too.
A rerun of the two settings on all eight took 15 minutes with 8 jobs, on a loaded 8-core machine,
and of all eleven settings 75 minutes with 6 jobs.
The documents aren't in the repository, and they appear here by kind only.
How to measure them again, or other documents, is in \texttt{measure-real/README.md}.

\subsection{Method}

Eight documents, all built before with the package on \textsf{memoir}: three sets of
lecture notes, one set of course notes, a book, two literate programs and a paper.
One is built with Xe\TeX, the others with pdf\TeX.
Together they have 814 pages and 249 lists.

Each document is copied once, with the generated files of its last build, and those are reused:
PythonTeX output, biber's and minted's.
All numbers are from that fixed snapshot, taken on 4 November 2026, so each setting and commit sees
the same text.
One set of lecture notes was woven and run through biber in its copy, as its woven source wasn't kept.
\texttt{audit.py} checks each copy before anything is counted, see below.
Then each document is built under each setting until its \texttt{.aux} doesn't change, at most six times.

A setting is a small change to a copy of \texttt{didactic.sty}:
\begin{description}
\item[current] the package as it is: club and widow penalties are \textsf{memoir}'s 1000, and all
  keep options have their defaults.
\item[1000/10000, 10000/10000] other club and widow penalties.
\item[keepleadin off] \texttt{keepleadin=false}.
\item[keep options off] also \texttt{keeplines=0}, \texttt{keepends=1},
  \texttt{keepafterheading=0} and \texttt{keepmaxgap=0}.
  The documents need the package's macros, so this is the nearest to plain \LaTeX\ with the
  package's fonts and margins.
\item[keepafterheading=1] that option's default changed.
\item[estimate of two lines] the estimate before a heading asks for two lines of room rather than one.
\item[list rule off] without the package's rule that keeps a list with its lead-in, so that the
  penalty before a list's first item is \textsf{memoir}'s $-51$.
\item[list 0, list 10000] the same, with that penalty changed for \texttt{itemize},
  \texttt{enumerate} and \texttt{description}.
\end{description}

A probe in the copy of the package writes, at shipout like a label, the page of each paragraph's
last line and of each list's first item, and for each page that goes through the main output
routine the penalty it broke at and its natural height.
A list is parted when those two pages differ.
A lone line is read from the penalty: 1000 or 1100 is a first line alone, and a separate setting
with a widow penalty of 1001 tells a last line from a first when both are 1000.
Its breaks are those of \emph{current}.
The empty space $e$ on a page is its height less its natural height, in baselines, so that a
negative $e$ means glue shrank.
It says nothing of where the space shows under \apcs{flushbottom}.

\subsection{Page counts}

\begin{center}\small
\begin{tabular}{@{}lrrrrrrrr@{}}
\toprule
& keep & keep- & 1000/1000 & 1000/ & 10000/ & keepafter- & list & list \\
Document & off & leadin off & (current) & 10000 & 10000 & heading=1 & rule off & 10000 \\
\midrule
Lecture notes & 121 & 121 & 121 & 121 & 121 & 121 & 121 & 121 \\
Lecture notes & 67 & 67 & 68 & 68 & 68 & 68 & 68 & 68 \\
Lecture notes & 108 & 108 & 108 & 108 & 108 & 108 & 108 & 108 \\
Course notes & 101 & 101 & 101 & 101 & 101 & 101 & 101 & 101 \\
Book & 124 & 124 & 124 & 124 & 124 & 124 & 124 & 124 \\
Literate program & 93 & 92 & 93 & 93 & 93 & 93 & 93 & 93 \\
Literate program & 77 & 77 & 78 & 78 & 78 & 78 & 78 & 78 \\
Paper & 120 & 121 & 121 & 121 & 122 & 121 & 121 & 121 \\
\midrule
Total & 811 & 811 & 814 & 814 & 815 & 814 & 814 & 814 \\
\bottomrule
\end{tabular}
\end{center}

So club 1000 with widow 1000 and with widow 10000 tie on these documents, and 10000 for both adds
a page to one of them.
The keep options cost three pages in 814, one each.
On the two literate programs the last page with them holds under ten words, so a small shift
upstream can change such a count either way.
With a list's penalty at 0 the total is also 814.

\subsection{Lone lines and empty space}

\begin{center}\small
\begin{tabular}{@{}lrr@{}}
\toprule
Club/widow & First line alone at the foot & Last line alone at the top \\
\midrule
1000/1000 (current) & 3 & 10 \\
1000/10000 & 4 & 0 \\
10000/10000 & 0 & 0 \\
\bottomrule
\end{tabular}
\end{center}

These are on the 676 pages that go through the main output routine, of 814; a page of floats isn't
counted.
A paragraph of two lines broken in the middle would count on both sides; none is.

\begin{center}\small
\begin{tabular}{@{}lrrr@{}}
\toprule
Setting & Pages $e\ge3$ & Pages $e>5$ & Largest $e$ \\
\midrule
keep options off & 104 & 62 & 48.6 \\
keepleadin off & 97 & 62 & 45.0 \\
current & 118 & 69 & 49.2 \\
1000/10000 & 120 & 69 & 49.2 \\
10000/10000 & 122 & 71 & 51.0 \\
keepafterheading=1 & 118 & 69 & 49.2 \\
list rule off & 118 & 70 & 49.2 \\
list 0 & 116 & 71 & 49.2 \\
list 10000 & 118 & 69 & 49.2 \\
\bottomrule
\end{tabular}
\end{center}

The largest values are the last pages of chapters and of documents, so compare the rows, not the
values.
Compared with the keep options off, \emph{current} has 14 more pages with $e\ge3$ and 7 more with
$e>5$, of 676.

\subsection{A list and its lead-in}

Of the 249 lists, 31 come right after a heading and aren't counted, which leaves 218.

\begin{center}\small
\begin{tabular}{@{}lrrrrrrr@{}}
\toprule
Document & Lists & After a heading & keep off & list rule off & list 0 & list 10000 & current \\
\midrule
Lecture notes & 5 & 0 & 0 & 0 & 0 & 0 & 0 \\
Lecture notes & 5 & 0 & 0 & 0 & 0 & 0 & 0 \\
Lecture notes & 15 & 0 & 1 & 1 & 1 & 0 & 0 \\
Course notes & 165 & 31 & 8 & 8 & 1 & 0 & 0 \\
Book & 19 & 0 & 1 & 1 & 1 & 0 & 0 \\
Literate program & 9 & 0 & 0 & 1 & 1 & 0 & 0 \\
Literate program & 2 & 0 & 0 & 0 & 0 & 0 & 0 \\
Paper & 29 & 0 & 4 & 3 & 3 & 0 & 0 \\
\midrule
Total & 249 & 31 & 14 & 14 & 7 & 0 & 0 \\
\bottomrule
\end{tabular}
\end{center}

The source was read for each of the 14 lists parted with the rule off.
Thirteen have a real lead-in at the foot of the page before, such as \enquote{Countermeasures
include:}, and three of those are a line right under a heading.
One is a list resumed after an explanation of the answer before it.
None follows a listing or table.
So the problem as asked is 13 lists in 218.
On each of those pages $e$ is at most 1.4 lines, so these pages are nearly full, and the empty space
is elsewhere.
In two cases a figure tops the next page, so that the list isn't its first thing.

What the rule costs is lead-ins going along:
\begin{center}\small
\begin{tabular}{@{}lrrrr@{}}
\toprule
& Lead-in goes & No longer parted, & Newly & Largest $e$ on \\
Setting & forward with the list & earlier pages shifted & parted & the page left \\
\midrule
list 0 & 3 & 5 & 1 & 2.8 \\
list 10000 & 8 & 6 & 0 & 3.9 \\
current & 8 & 6 & 0 & 3.9 \\
\bottomrule
\end{tabular}
\end{center}

On these documents \emph{current} has the same words on each page as \emph{list 10000}.
They differ where the rule is released past \texttt{keepmaxgap}, which doesn't happen to a list here.

\subsection{A heading}

With \texttt{keepafterheading=1}, instead of 2, every document has the same words on every page, and
the same 814 pages.
That isn't a setting lost on the way: a trace in the build of the lecture notes and of the course
notes, where the words changed at the earlier commit below, shows 1 or 2 reaching
\apcs{@didactic@keep@room} and the lines owed in the first paragraph.
On a generated heading the setting does change the page, when the first paragraph has one line:

\begin{center}\small
\begin{tabular}{@{}llll@{}}
\toprule
Under the heading & Room in lines & \texttt{keepafterheading=1} & \texttt{keepafterheading=2} \\
\midrule
A paragraph of many lines & 4--6 & the same & the same \\
A list of one-line items & 4--6 & the same & the same \\
One line, then a paragraph & 4--4.5 & stays with that line & moves \\
One line, then a paragraph & 5 & stays with that line & stays with two lines \\
One line, then a paragraph & 5.5--6 & the same & the same \\
\bottomrule
\end{tabular}
\end{center}

So the setting decides only what the chain holds after a first paragraph of one line, since \LaTeX\
already keeps two lines of a longer one, and a list under a heading has what \LaTeX\ gives it.
None of these documents has a heading like that at the foot of a page.

At commit \texttt{dd90b97}, before the rule for a list under a heading, the setting changed the words
on 10 pages, with 118 and 70 pages of $e\ge3$ and $e>5$ against 114 and 68.
There the estimate before a heading asked for \texttt{keepafterheading} lines of room, and now it
asks for one, as it can't know what follows.
The row with an estimate of two lines shows what that change did at the current code: 814 pages,
with words differing on 4 pages (3 in one set of lecture notes, 1 in the course notes), and 116 and 68 pages of $e\ge3$ and $e>5$ against 118 and 69.
On one of them, the heading of an exercise with two lines at the foot of a page moves to the next page.

A list directly under a heading was measured at \texttt{dd90b97} with a scratch change that let its
first item stay with one line, against the two-line rule: no page differed in words.
That is now the package's rule.

\subsection{Warnings}

Without the probe, \emph{current} gives 19 warnings on all eight documents: 12 about a heading, 6
that a tie is released, and 1 other.
With the keep options off there are none, and with \texttt{keepleadin} off 12, all about headings, both with the probe.
With the list rule off there are 25 against 21, both with the probe.

\subsection{The pair of penalties on generated text}

The real documents tie, so the choice between the pairs rests on generated text.
These are \texttt{measure-flow/} at the current code, with \textsf{noweb} or \textsf{minted} chunks,
and \texttt{queries/pen.py}: 6 seeds of 300 paragraphs in \textsf{memoir} and \textsf{article}, with
and without headings, with no package, which doesn't depend on the commit.

\begin{center}\small
\begin{tabular}{@{}lrrr@{}}
\toprule
Pages & 1000/1000 & 1000/10000 & 10000/10000 \\
\midrule
Flows, \textsf{noweb} chunks & 210 & 210 & 210 \\
Flows, \textsf{minted} chunks & 194 & 195 & 195 \\
\texttt{pen.py}, \textsf{memoir} & 613 & 616 & 617 \\
\texttt{pen.py}, \textsf{article} & 610 & 610 & 610 \\
\midrule
Total & 1627 & 1631 & 1632 \\
\bottomrule
\end{tabular}
\end{center}

So 1000/1000 is 4 pages shorter than 1000/10000 in about 1630.
Three of those 4 are in \texttt{pen.py} in \textsf{memoir}, and one is a flow whose last page holds
about a line.
The price is in lone lines:

\begin{center}\small
\begin{tabular}{@{}lrrrrrr@{}}
\toprule
& \multicolumn{2}{c}{1000/1000} & \multicolumn{2}{c}{1000/10000} & \multicolumn{2}{c}{10000/10000} \\
Lone lines & first & last & first & last & first & last \\
\midrule
Flows, \textsf{noweb} chunks & 2 & 3 & 1 & 0 & 0 & 0 \\
Flows, \textsf{minted} chunks & 4 & 7 & 2 & 0 & 0 & 0 \\
\texttt{pen.py} & 55 & 59 & 63 & 0 & 0 & 0 \\
\bottomrule
\end{tabular}
\end{center}

\subsection{Generated literate programs}

\texttt{measure-flow/} generates 12 documents of 60 pairs of a code chunk and text each, in \textsf{memoir} with \texttt{article} and \texttt{oneside}, with the lengths of chunks and the share of a lone \texttt{@} measured on the
maintainer's \texttt{.nw} files.
They have no margin notes, floats or displays.
The chunks are set by \textsf{noweb} or by \textsf{minted} as the \texttt{tominted} filter weaves them,
with filler code.
Plain means no package, with its fonts and margins.

\begin{center}\small
\begin{tabular}{@{}llrrrrrrr@{}}
\toprule
Chunks & Setting & Pages & Mean $e$ & Largest $e$ & $e\ge2$ & $e\ge3$ & $e\ge5$ & Warnings \\
\midrule
noweb & plain, no option & 226 & 7.92 & 41.1 & 154 & 136 & 101 & -- \\
noweb & plain, breakcode & 210 & 0.84 & 6.2 & 48 & 31 & 2 & -- \\
noweb & package, no option & 225 & 7.81 & 41.1 & 154 & 136 & 100 & 14 \\
noweb & package, breakcode & 210 & 0.65 & 5.6 & 33 & 20 & 1 & 1 \\
noweb & package, breakcode, keepmaxgap=0 & 211 & 0.95 & 7.6 & 50 & 33 & 4 & 0 \\
minted & plain, breakcode & 196 & 0.96 & 6.2 & 45 & 28 & 4 & -- \\
minted & package, breakcode & 194 & 0.71 & 4.5 & 34 & 17 & 0 & 6 \\
minted & keepmaxgap=0 & 197 & 1.48 & 8.9 & 64 & 43 & 17 & 0 \\
minted & keepleadin=false & 194 & 0.59 & 5.2 & 33 & 12 & 1 & 2 \\
minted & keeplines=0 & 194 & 0.64 & 4.5 & 32 & 15 & 0 & 6 \\
\bottomrule
\end{tabular}
\end{center}

Without \textsf{noweb}'s \texttt{breakcode}, its \apcs{filbreak} leaves up to 41 lines empty on a page,
and the package doesn't touch that break.
With it the package leaves fewer short pages than plain.
\textsf{noweb}'s own \enquote{leave no more than this}, 0.2\,in, isn't a bound: plain, a page ends
6.2 lines short.
The one page over \texttt{keepmaxgap} with \textsf{noweb} chunks ends after a chunk, at
\textsf{noweb}'s own penalty.

The package drops \textsf{noweb}'s \apcs{nwdocspar} under \texttt{keepmaxgap}.
With \textsf{noweb} chunks and \texttt{breakcode}, in 210 pages, which have 198 breaks since a
document's last page ends by itself, that takes the breaks at
$-200$ from 37 to 0, those after a chunk at $-301$ from 60 to 76, and those inside a chunk at 301 from
93 to 112: 19 more, against the smaller gaps in the table.

\texttt{keepmaxgap=0} also keeps \apcs{nwdocspar}, so the two were measured apart: the limit on with
the break kept, by taking the \apcs{let} out of a copy of the package, and the limit off with the
break dropped, by a \apcs{let} in the preamble.

\begin{center}\small
\begin{tabular}{@{}lllrrrrr@{}}
\toprule
Chunks & \texttt{keepmaxgap} & \apcs{nwdocspar} & Pages & Mean $e$ & Largest $e$ & $e\ge3$ & $e\ge5$ \\
\midrule
noweb & 5 & dropped (current) & 210 & 0.65 & 5.6 & 20 & 1 \\
noweb & 5 & kept & 211 & 0.95 & 7.6 & 33 & 4 \\
noweb & 0 & dropped & 210 & 0.65 & 5.6 & 21 & 1 \\
noweb & 0 & kept & 211 & 0.95 & 7.6 & 33 & 4 \\
minted & 5 & dropped (current) & 194 & 0.71 & 4.5 & 17 & 0 \\
minted & 5 & kept & 196 & 1.01 & 6.2 & 28 & 3 \\
minted & 0 & dropped & 196 & 1.44 & 10.5 & 39 & 14 \\
minted & 0 & kept & 197 & 1.48 & 8.9 & 43 & 17 \\
\bottomrule
\end{tabular}
\end{center}

With \textsf{noweb} chunks all of 7.6 against 5.6 is \apcs{nwdocspar}, and the limit changes one
page's count in $e\ge3$, as those chunks can break anywhere at a moderate penalty.
With \textsf{minted} chunks, which the package keeps together, the limit does the larger part, and
dropping the break matters only along with it.

\subsection{Build time}

One \LaTeX\ pass at v2.13 (\texttt{743a025}), before the keep options, against the current code, with
each \texttt{.aux} settled first and three passes in turn for each version, by
\texttt{measure-real/timing.py}.
CPU seconds are the mean of the three, with their range.

\begin{center}\small
\begin{tabular}{@{}lrrrrrr@{}}
\toprule
& \multicolumn{2}{c}{Pages} & \multicolumn{2}{c}{\texttt{.aux}, kB} & \multicolumn{2}{c}{CPU s} \\
Document & v2.13 & current & v2.13 & current & v2.13 & current \\
\midrule
Lecture notes & 67 & 68 & 150 & 179 & 31.8 (26.6--37.8) & 33.3 (24.0--38.6) \\
Literate program & 77 & 78 & 45 & 63 & 13.8 (11.5--15.2) & 13.9 (12.2--17.2) \\
Paper & 121 & 122 & 156 & 188 & 35.6 (32.7--38.6) & 35.0 (26.2--39.6) \\
Generated, 3000 sections & 250 & 250 & 318 & 710 & 4.9 (4.2--5.3) & 5.8 (5.6--6.0) \\
\bottomrule
\end{tabular}
\end{center}

On the real documents no difference shows above the noise.
Other jobs held the load at 8 to 17 on an 8-core machine, and passes of one version vary by up to half,
so a real difference under about 20 per cent on these documents wouldn't show.
The generated document, with a section every few lines, has an \texttt{.aux} over twice the size and is
about a fifth slower here, not twice.
The paper's 122 pages come from settling with only four passes.

\subsection{Checks a new measurement must do}

Each of these went wrong once, without showing in the page count.
\begin{itemize}
\item An \apcs{input}\{contents.tex\} with no woven file found a test fixture of that name
  through \texttt{TEXINPUTS}: a lecture counted as a shorter document.
  \texttt{audit.py} fails on a file read from outside the copy, by \texttt{-recorder}.
\item Reused PythonTeX output named files by absolute path, read through \textsf{minted}, which
  \texttt{-recorder} doesn't list.
  The audit also looks for absolute paths in reused output, and \texttt{prepare.py} points them at the copy.
\item A probe that wrote a whatsit into an empty paragraph, such as \textsf{noweb}'s
  \apcs{noindent}\apcs{par}, made a line of it and added pages.
  Compare with the setting without the probe.
\item \texttt{biber} with a wrong path emptied the \texttt{.bbl}, which shows only as fewer pages.
  \texttt{prepare.py} fails on an empty one.
\item A quota reached on the scratch disk, with \texttt{df} showing free space, can leave a short log.
  A build with a quota line or without its output line is no result, and \texttt{analyse.py} refuses it.
\item Builds killed by a timeout left their \LaTeX\ running, which wrote into the next build of that directory.
\item Delete a build only after its numbers are saved.
\end{itemize}

\subsection{Limits}

\begin{itemize}
\item Eight documents, one author, and one set of course notes has 165 of the 249 lists and 8 of the
  14 parted ones.
\item Generated files are as last built, the same for every setting.
  Some woven files are older than their sources, by up to 25 days, and one document has 9 undefined
  citations from an older bibliography, in every setting.
\item The probe still moves some lines between pages in two documents, on 13 and 22 pages, with the
  same page count.
  The page counts and the warnings for \emph{current} are from a build without it, and so is the
  check of the list rule off in those two documents.
  List counts come from the builds with the probe.
  In those two documents the parted lists were also read in the text of a build without it, and each
  falls on the same pages.
\item A list is measured against the last paragraph before it, whatever lies between: hence the reading above.
\item The choice of breaks doesn't depend on a flush or ragged bottom, as \apcs{flushbottom} and
  \apcs{raggedbottom} act after a page is broken.
\item A rerun on documents that have moved on since the snapshot mixes their changes with the
  package's, so keep a snapshot when comparing commits.
  One literate program here changed on the day of the snapshot: from its source it now has 79 pages,
  from the snapshot 78.
\end{itemize}
